% TOP_PROBLEM: {"id": 2905, "title": "Kirby Problem 4.29", "category_id": 7, "category": {"id": 7, "name": "topology", "display_name": "Topology", "description": "Properties preserved under continuous deformations.", "slug": "topology", "order_index": 7, "created_at": "2026-07-31T15:26:25.671Z"}, "status": "open"}
% FIRST_POSTED: null
% ENTRY_KIND: statement_only
% Imported from: batch13.tex; UnsolvedMath ID 2905
\documentclass[11pt]{article}
\usepackage[margin=25mm]{geometry}
\usepackage{fontspec}
\setmainfont{Latin Modern Roman}
\usepackage{amsmath,amssymb,mathtools}
\usepackage{unicode-math}
\setmathfont{Latin Modern Math}
\usepackage{xurl,hyperref}
\hypersetup{colorlinks=true,urlcolor=blue}
\setcounter{secnumdepth}{0}
\setlength{\parindent}{0pt}
\setlength{\parskip}{5pt}
\setlength{\emergencystretch}{3em}
\newcommand{\sref}[2]{\hyperlink{source:#1:#2}{[S#2]}}
\newcommand{\cataloguescope}[1]{\paragraph{Recorded target.}#1}
\begin{document}
% UnsolvedMath ID 2905; source catalogue code KP-4.29
\setcounter{section}{2904}
\section{Kirby Problem 4.29}\label{2905}
{\small\sffamily UnsolvedMath ID 2905 · Topology}
\subsection{Definitions and mathematical statement}

\paragraph{Shared notation.}$\mathbb N=\{1,2,\ldots\}$, and $\mathbb Z,\mathbb Q,\mathbb R,\mathbb C$ denote the integers, rationals, reals and complex numbers. Any other conventions are specified locally.


A locally flat embedded surface $\Sigma\subset S^4$ is a closed connected two-manifold with neighborhoods modeled on the inclusion $\mathbb R^2\subset\mathbb R^4$. Suppose its complement has cyclic fundamental group, meaning a group generated by a single element (finite or infinite).

For orientable genus $g\ge0$, the standard unknotted surface is the boundary of a standard genus-$g$ handlebody in an equatorial $S^3\subset S^4$. For the nonorientable models, write $P_+$ and $P_-$ for the standard unknotted projective planes with normal Euler numbers $+2$ and $-2$. One explicit model is the normalized Veronese embedding
\[
\mathbb{RP}^2\longrightarrow S^4\subset\operatorname{Sym}_0(3,\mathbb R),\qquad
[v]\longmapsto\sqrt{\frac32}\left(vv^{\mathsf T}-\frac13I_3\right),\quad |v|=1,
\]
where $\operatorname{Sym}_0(3,\mathbb R)$ is the five-dimensional space of real symmetric trace-zero matrices with its Euclidean matrix norm; reflection gives the other sign. The standard nonorientable genus-$h$ surfaces are connected sums of $h$ copies of $P_+$ or $P_-$, formed by joining standard local disks by tubes. Their normal Euler number is the sum of the signs $\pm2$. An embedding is topologically or smoothly unknotted if it is ambiently isotopic, in that category, to a corresponding standard model.
\paragraph{Statement recorded in the supplied source.}
Answer both parts:
\begin{enumerate}
\item[(a)] Must every locally flat $\Sigma\subset S^4$ with $\pi_1(S^4\setminus\Sigma)$ cyclic be topologically unknotted?
\item[(b)] If $\Sigma$ is smoothly embedded and orientable, must it be smoothly unknotted, or is there a counterexample with cyclic complement group?
\end{enumerate}
Part (a) includes nonorientable surfaces of every genus and normal Euler number that can occur; part (b) is restricted to orientable surfaces. These universal targets are not replaced by the particular cases still unsettled in the source remarks. The source also records:
\begin{enumerate}
\item[(i)] If $\Sigma$ is smooth and a standard height (radial-level) function $h:S^4\to\mathbb R$ restricts to a Morse function on $\Sigma$ with exactly one local minimum, must $\Sigma$ be smoothly unknotted? Must it be topologically unknotted? Here the standard levels of $h$ are three-spheres between its two poles.
\item[(ii)] If $\Sigma=F_+\cup_U F_-$ is the union of two ribbon surfaces in the opposite four-ball hemispheres along an unknotted circle $U\subset S^3$, must $\Sigma$ be smoothly unknotted? Must it be topologically unknotted? A ribbon surface is a proper smooth surface for which the radial function in its four-ball has no interior local maxima; smooth the joined surface along $U$.
\end{enumerate}
No cyclic-complement assumption is added to these two separately recorded questions. \sref{2905}{2}
\subsection{Short English statement}
Ask whether cyclic complement group forces topological or smooth unknotting, and whether a single minimum or a decomposition into two ribbon surfaces with unknotted cross-section forces unknotting in either category.
\subsection{Sources}
\begin{description}
\item[\hypertarget{source:2905:1}{S1}] ULAM.AI and the UnsolvedMath Contributors, UnsolvedMath: A Curated Collection of Open Mathematics Problems, record 2905 (KP-4.29). CC BY 4.0. \url{https://huggingface.co/datasets/ulamai/UnsolvedMath}.
\item[\hypertarget{source:2905:2}{S2}] R. İnanç Baykur, Robion C. Kirby and Daniel Ruberman, K3—A New Problem List in Low-Dimensional Topology, Mathematical Surveys and Monographs 295 (2026), author version, Problem 4.29 and Remarks, PDF page 213. \url{https://math.berkeley.edu/sites/default/files/surv-295-ruberman-watermarked-author-pdf.pdf}.
\end{description}
\label{end:2905}
\end{document}
